\chapter{CONCLUSION}

This project set out to detect cross-bank mule-account laundering --- a crime that hides in the gaps between banks --- and to measure how well such a detector holds up when criminals adapt. ARGUS represents the flow of money as a heterogeneous, temporal graph of accounts, owning entities and banks, and detects laundering with an edge-level graph neural network whose neighbourhood sampling never looks into the future.

On the IBM AML HI-Small benchmark, evaluated on the chronologically last fifth of 5.08 million transactions, the GNN reached a PR-AUC of 0.726 and an F1 of 0.729 for the laundering class, against 0.296 and 0.367 for LightGBM on the same features and 0.002 and 0.006 for fixed-threshold rules. It fully recovered 60.4\% of laundering rings, against 4.2\% for LightGBM, confirming that the network view is what makes ring-level detection possible. Removing look-ahead from the sampler changed PR-AUC by less than 0.01, so the gain is not a leakage artefact.

The central contribution is the adversarial evaluation. A multi-bank digital twin built on real background traffic, with injected rings of fully known structure, allowed six attack tactics to be measured individually and in an adaptive multi-round game. The GNN largely neutralised structuring (degradation 0.038 F1 against 0.439 for LightGBM) and lost a third as much as LightGBM to feature mimicry; under an attacker who adapted four times it retained a retrained F1 of 0.878 against 0.834. Relation camouflage, in which mules surround themselves with legitimate counterparties, emerged as the open problem: it removed about 0.44--0.48 F1 from both models, and retraining recovered it only partly.

Finally, ARGUS turns detections into decisions a bank can act on. The investigation agent reconstructed detected rings with a Jaccard similarity of 0.90 and identified the victim in 95\% of cases on clean data; it drafts Suspicious Transaction Reports whose language-model narrative is verified number by number against computed facts, and it leaves every freeze decision to a human officer. A five-step live demonstration presents the alert, the network, the investigation and the attack.

In short, ARGUS shows that graph intelligence detects cross-bank laundering far better than per-transaction models, that it is substantially more robust to the most common evasion tactics, and --- by measuring degradation and recovery rather than a single accuracy figure --- it identifies relation camouflage as the attack that the next generation of AML detectors must be designed against.
